\section{10 декабря 1926 г.}

\lp{20260606_184605782}

Дорогой мой и милый Коля.

Вчера, 9 декабря, я получил твою скрипичную сонату. Сердечное тебе спасибо за
твой подарок и за посвящение. Я вчера же вечером с ней ознакомился и с
наслаждением несколько раз проиграл ее.
Сочинение это меня глубоко тронуло и захватило с первых же нот. Изложена соната
мастерски и я жажду поскорее услышать ее со скрипкой, что и сделаю на днях.

\rem{З...} также получил твою сонату в тот же день, что и я, и очень тебя
благодарит. Писать ему некогда, так как он наспех переинструментовывает оперетту
Лекока для малого оркестра камерного театра.

\lp{20260606_184559271}

Время летит стрелой и уже недолго осталось до твоего приезда в Москву. С
нетерпением жду свидеться с тобой и поглядеть на тебя. -- Музыкальная жизнь в
Москве кипит ключём, но вода в этом ключе не всегда бывает пригодна для питья.
За последнее время, чаще, чем в прежние годы, я стал прихварывать гастритами и
мне теперь иногда кажется, что это явление может быть находится в некоторой
связи с частым посещением концертов.

Недавно я играл в Персимфансе свой новый органный концерт со струнным
оркестром, который написал летом. Сыграли его очень недурно (в техническом
отношении вполне удовлетворительно). В этом же

\lp{20260606_184559271}

\noindent
концерте играли вступление к моей опере. -- Живу я (если вообще можно назвать
жизнью суету и суматоху) терпимо, но без отдыха и без досуга для своей работы,
что конечно раздражает, утомляет и отравляет настроение.

Надеюсь, что ко дню твоего рождения я соберусь написать тебе еще письмо, но на
всякий случай поздравляю тебя с наступающим днем рождения, крепко тебя целую и
от души желаю здоровья и возможности так обставить свою жизнь, чтобы как можно
меньше было помех для твоей композиторской работы и как можно больше опусов
таких как твоя 2-я скрипичная соната и песни op. 45 (до сих пор в Москве нет
op. 43 и 46).

\lp{20260606_184605782}

Напиши мне два слова, если найдешь минутку, не будешь ли ты иметь что-либо
против, если я дам сыграть твою 2-ю скрипичную сонату в одном из ученических
концертов хорошим ученикам моего камерного класса (у меня имеются там отличные
скрипачи и очень хорошие пианисты). Без твоего согласия мне бы этого делать не
хотелось.

Передай от меня и Екатерины Петровны самый сердечный поклон Анне Михайловне.

Еще раз крепко тебя целую. Твой А. Гедике.

Если будешь писать Миле и С. В. \rem{Рахманинову} передай им мой сердечный
поклон. Поцелуй от меня Льва Эдуардовича и поклонись Ольге Николаевне.

19 10/XII 26 г.
